\documentclass[12pt,a4paper]{article}
\usepackage[margin=3cm]{geometry}
\usepackage{graphicx}
\usepackage{times}
\usepackage[hidelinks]{hyperref}

% ==== ISI TEMPLATE DI SINI ====
\newcommand{\judulLaporan}{LAPORAN [MATA KULIAH]}
\newcommand{\subJudul}{``[JUDUL / TEMA LAPORAN]''}
\newcommand{\logoFile}{logo.png}
\newcommand{\dosenPengampu}{[Nama Dosen 1]\\{}[Nama Dosen 2]}
\newcommand{\penyusun}{%
[Nama Mahasiswa 1] & ([NIM]) \\{}
[Nama Mahasiswa 2] & ([NIM]) \\{}
[Nama Mahasiswa 3] & ([NIM]) \\{}
}
\newcommand{\prodi}{PRODI [Nama Prodi]}
\newcommand{\jurusan}{JURUSAN [Nama Jurusan]}
\newcommand{\fakultas}{FAKULTAS [Nama Fakultas]}
\newcommand{\universitas}{UNIVERSITAS [Nama Universitas]}
\newcommand{\tahun}{2026}
% ==== AKHIR BAGIAN YANG PERLU DIUBAH ====

\begin{document}

\begin{titlepage}
\centering
\vspace*{1cm}
{\bfseries \judulLaporan \par}
\vspace{0.8cm}
{\bfseries \subJudul \par}
\vfill
\includegraphics[width=4.5cm]{\logoFile}
\vfill
{\bfseries DOSEN PENGAMPU :\par}
\vspace{0.2cm}
{\bfseries \dosenPengampu \par}
\vspace{1cm}
{\bfseries DISUSUN OLEH :\par}
\vspace{0.3cm}
\begin{tabular}{l@{\hspace{2cm}}l}
\penyusun
\end{tabular}
\vfill
{\bfseries \prodi \par}
\vspace{0.3cm}
{\bfseries \jurusan \par}
\vspace{0.3cm}
{\bfseries \fakultas \par}
\vspace{0.3cm}
{\bfseries \universitas \par}
\vspace{0.6cm}
{\bfseries \tahun \par}
\end{titlepage}

\tableofcontents
\newpage

\section*{Kata Pengantar}
[Isi kata pengantar]

\newpage
\section{Pendahuluan}
\subsection{Latar Belakang}
[Isi latar belakang]

\subsection{Rumusan Masalah}
[Isi rumusan masalah]

\subsection{Tujuan}
[Isi tujuan]

\section{Pembahasan}
[Isi pembahasan, contoh sitasi: \cite{lankton2008introduction}]

\section{Penutup}
\subsection{Kesimpulan}
[Isi kesimpulan]

\subsection{Saran}
[Isi saran]

\bibliographystyle{apalike}
\bibliography{pustaka}

\end{document}
